\documentclass[11pt]{article}

\usepackage[margin=0.82in]{geometry}
\usepackage{amsmath,amssymb,booktabs,array,longtable,enumitem,microtype,xcolor,fancyhdr,lastpage}
\usepackage[T1]{fontenc}
\usepackage{lmodern}
\usepackage{titlesec}
\usepackage{parskip}
\usepackage{caption}
\usepackage{hyperref}

\hypersetup{colorlinks=true,linkcolor=black,urlcolor=blue,citecolor=black}

\definecolor{darkblue}{RGB}{26,54,93}

\titleformat{\section}{\large\bfseries\color{darkblue}}{\thesection}{0.6em}{}
\titleformat{\subsection}{\normalsize\bfseries}{\thesubsection}{0.6em}{}
\titleformat{\subsubsection}{\normalsize\bfseries}{\thesubsubsection}{0.6em}{}

\setlist[itemize]{leftmargin=1.3em,itemsep=0.2em,topsep=0.25em}
\setlist[enumerate]{leftmargin=1.5em,itemsep=0.2em,topsep=0.25em}
\setlength{\tabcolsep}{3pt}
\setlength{\LTleft}{0pt}
\setlength{\LTright}{0pt}

\pagestyle{fancy}
\fancyhf{}
\lhead{Gamma Tracker V6}
\rhead{Implemented measurement}
\cfoot{\thepage\ / \pageref{LastPage}}
\setlength{\headheight}{14pt}

\newcommand{\code}[1]{\texttt{\detokenize{#1}}}
\newcommand{\eps}{\varepsilon}

\title{\vspace{-1.5cm}\textbf{Gamma Tracker V6}\\[0.25em]
\large What the live page measures, how each chart is drawn, and how
to read the surface}
\author{}
\date{Implementation as of 22 September 2026}

\begin{document}
\maketitle
\thispagestyle{fancy}
\vspace{-1.2em}

\begin{abstract}
This note describes only Gamma Tracker V6
(\code{Z:/uikitxv2/OptMonitor/gamma_screener/gamma_ladder_standalone/gamma_tracker_dash_v6.py},
port 8055, local \url{http://127.0.0.1:8055/},
Tailscale \url{https://aasthajain.tail174727.ts.net:8443/}).
The page is a live measurement of the options implied-vol
surface against its own recent history. It shows ATM backbone, risk-reversals, butterflies,
a percentile heatmap of the current smile versus sampled history, and
an event-conditioned view of those same ranks.
\end{abstract}

\tableofcontents
\newpage

\section{Goal}

Where is the live option surface
relative to its own recent distribution, and is that location ordinary
for the current time of day, vol regime, and event window?

The table, heatmap, timeseries, and smile are different views of that
same construction:

\begin{itemize}
    \item The summary table is the latest snapshot: ATM vol, weighted
    skew, weighted curvature, and the raw RR / fly legs.
    \item The heatmap ranks each delta-bucket versus sampled history
    (0--100).  Red is high in that history; blue is low.
    \item The metric graph is the same numbers over time.
    \item The smile is the latest snapshot in strike space, not
    percentile space.
    \item Event scores do not change the live IVs.  They change
    \emph{which historical observations} the ranks are compared to.
\end{itemize}

A cell at the 90th percentile means ``this wing is high versus the
chosen history bucket.''  It is not a forecast that vol will fall.


\section{Page, stack, and data}


\subsection{Sources}

Live option rows come from AccessAPI
\code{/live/options/snapshot} with columns
\code{ts_event}, \code{symbol}, \code{strike_price}, \code{expiration},
\code{futures_mid}, \code{iv_mid}, \code{delta_F_mid},
\code{gamma_F_mid}, and related greeks.
The adapter polls that endpoint every 1 second and treats the feed as
disconnected if it is older than 10 seconds.

Futures mid is the median of \code{futures_mid} on the option rows.
If that is missing or non-positive, the writer falls back to
AccessAPI TY TBBO mid
(\code{/live/tbbo/recent?symbols=TY\&n=1}).
A snapshot is skipped until a positive futures price exists.

Event-conditioned percentiles still read
\code{Y:/Aastha/market_data_puller/workspace/economic_events.db}
(\code{bbg_data} tiers 1--3 plus selected \code{economic_events}
rows).  The upcoming-events banner reads Yaman
\code{Y:/Yaman/Fundamentals/Economic Events/Data/events.sqlite}
instead; the old calendar currently has no rows after 2026-08-14.
Event weights live in
\code{event_params.json} (calibrated 2026-03-25).

The writer does not use CMED Excel or Redis.

\subsection{Database}

Path: \code{C:/uikitxv2_cache/gamma_tracker/gamma_tracker_v3.db}.

\begin{center}
\begin{tabular}{@{}lp{0.62\linewidth}@{}}
\toprule
Table & Role \\
\midrule
\code{snapshot_v2} & One row per persist: timestamp, source key,
\code{file_mtime}. \\
\code{summary_metrics_v2} & Per-expiry backbone, RR, fly, indices,
ROC. \\
\code{option_surface_v2} & Per-strike live surface for that snapshot. \\
\code{heatmap_fact} & Thinned history used to build percentiles.
Never pruned. \\
\code{futures_vol_log} & Futures mid plus the two front ATM vols.
Not drawn on V6. \\
\bottomrule
\end{tabular}
\end{center}

Raw snapshot / surface / summary / futures-vol rows older than 30 days
are deleted.  \code{heatmap_fact} is kept so the percentile cache can
span more than the raw-snapshot window.

A sidecar \code{reference_surface.json} stores the 4-hour median of
each RR and fly by expiry.  V6 does not read it.

A portable copy of the same page lives in
\code{Z:/uikitxv2/OptMonitor/gamma_screener/gamma_tracker_v6_standalone}
(writer, dashboard, slim \code{gamma_tracker_v6.db}).
It does not import the v2 persist library or the v3 writer.
The live page on port 8055 still reads
\code{gamma_tracker_v3.db} and was not switched.


\section{How a quote becomes a surface row}

\subsection{Valid instruments}

The adapter keeps a row if it has an expiration, instrument class
\(\in\{C,P\}\), and a numeric strike.  It then requires numeric
\(\delta=\code{delta_F_mid}\) and \(\sigma=\code{iv_mid}\).

Implied vol is stored as ``daily BP.''  

\subsection{Expiry ordinals}

The adapter first maps future weekday expiries, Chicago date, skipping
today and weekends, onto \code{1st}, \code{2nd}, \ldots.
The v3 writer then remaps again: every remaining expiry date strictly
after today is sorted and assigned
\code{2nd}, \code{3rd}, \ldots, \code{10th}.
The front remaining expiry is therefore labeled \code{2nd}.
\code{1st} is never persisted on this path.

Weekend expiries never enter the adapter map, so they never reach the
writer.

\subsection{Moneyness buckets}

Moneyness is the distance between the option strike \(K\) and the
current TY futures price \(F\), in hundredths of a point, rounded
to the nearest of \(25\), \(50\), \(75\), \(100\) (and \(0\) when
the strike is at the money):

\begin{equation}
d=100\cdot|K-F|,\qquad
m
=
\text{nearest of }\{0,25,50,75,100\}\text{ to }d.
\end{equation}

Example: if live TY is \(106\), then \(K=106.25\) is \(25\) out
and \(K=107\) is \(100\) out.  Bucket 0 is the ATM column.
After binning, each \((\mathrm{expiry},\mathrm{type},m)\) keeps one
strike.  Ties break to the lower strike.  Extra prints in the same
bucket are dropped.

\subsection{ATM vol on an expiry}

For each expiry date, ATM vol is the IV of the option---call or
put---whose strike is closest to the current TY futures price
(\(m=0\)).  That single number is copied onto every surface row
for that expiry as \code{Atm Vol (Daily BP)}.


\subsection{Cadence}

The writer persists only when all of these hold: the adapter is
connected, at least 30 seconds have passed, the live stream id
changed, futures price is positive, and
\code{compute_summary_metrics} produced at least one expiry row.
If the stream id has not changed, the same snapshot stays current
for longer than 30 seconds.

The top summary table is a separate clock.  The page re-reads the
latest \code{summary_metrics_v2} rows and repaints that table every
2 seconds.  Backbone, RR, fly, the two indices, and the 5/15/30-minute
ROC cells change only when a new snapshot is persisted.  The 6-hour
deltas and the status-line \code{ts_slope} are recomputed on each
2-second paint from snapshots already stored.

Heatmap-fact rows are written on every 60th successful snapshot
(about 30 minutes when the loop is clean).
Those facts are what the percentile cache later ranks against.
The heatmap colors follow the 15-minute cache rebuild, not the
2-second table paint.


\section{The summary metrics}

All of the following are computed per expiry, on the deduped
delta-bucket snapshot, in daily-BP units.

Let \(\sigma_m^C\) and \(\sigma_m^P\) be the call and put IVs in
bucket \(m\), and let \(\sigma^{\mathrm{ATM}}\) be that expiry's ATM
vol.

\subsection{Backbone}

\begin{equation}
B
=
\sigma^{\mathrm{ATM}}.
\end{equation}

That is the number in the table column \code{backbone} and the
left-hand ATM story on the page.

A high \(B\) means ATM options on that expiry are quoting more vol:
the market is charging more for a given near-spot move.  

High versus low is only meaningful against something: this expiry
earlier today (\code{backbone_roc_*} and the 6h delta), the other
expiries on the same snapshot (\code{ts_slope}), or the heatmap ATM
percentile.  A 5.2\,bp backbone that is ordinary for CPI week is
not the same story as 5.2\,bp on an empty calendar

\subsection{Risk reversals}

\begin{equation}
\mathrm{RR}_m
=
\sigma_m^C-\sigma_m^P,
\qquad
m\in\{25,50,75,100\}.
\end{equation}

Positive RR means the call wing is richer than the matched put wing.
Negative RR is put-skew bid.  Missing either leg leaves that RR
blank; it does not invent a value.

\subsection{Butterflies}

\begin{equation}
\mathrm{Fly}_m
=
\tfrac12\bigl(\sigma_m^C+\sigma_m^P\bigr)
-
\sigma^{\mathrm{ATM}}.
\end{equation}

Positive fly means the wings are rich versus ATM (smile is curved
up).  Negative fly is a flat or inverted smile versus ATM.

\subsection{Weighted indices}

Raw weights before normalization are
\(w_{25}=0.75\), \(w_{50}=0.50\), \(w_{75}=0.25\),
\(w_{100}=0.10\).  They are then divided by their sum \(1.6\):

\begin{equation}
\hat w_{25}=0.46875,\quad
\hat w_{50}=0.3125,\quad
\hat w_{75}=0.15625,\quad
\hat w_{100}=0.0625.
\end{equation}

\begin{align}
\mathrm{skew\_index}
&=
\sum_m \hat w_m\,\mathrm{RR}_m,\\
\mathrm{curve\_index}
&=
\sum_m \hat w_m\,\mathrm{Fly}_m.
\end{align}

A missing RR or fly at \(m\) contributes nothing to that index for
that term; the remaining weights are not re-normalized on the fly.

\subsection{Backbone rate of change}

Only expiries \code{2nd} through \code{5th} get ROC.
For window \(W\in\{5,15,30\}\) minutes, the writer finds the latest
earlier snapshot with \code{file_mtime} \(\le t-W\) and sets

\begin{equation}
\mathrm{ROC}_W
=
B_t-B_{t-W}.
\end{equation}

If that earlier snapshot has no backbone for the expiry, ROC is
blank.  This is a level change in daily BP, not a percent change.

\subsection{Term-structure slope}

Not stored.  The dashboard computes it from the latest table:

\begin{equation}
\mathrm{ts\_slope}
=
B_{2\mathrm{nd}}-B_{5\mathrm{th}}.
\end{equation}

Positive means the front ATM is richer than the fifth expiry.

\subsection{Six-hour deltas on the table}

If a snapshot exists at or before \(t-6\mathrm{h}\), the table adds
\code{backbone_delta_6h}, \code{skew_index_delta_6h}, and
\code{curve_index_delta_6h} as current minus that earlier snapshot.

\subsection{\code{backbone_pct}}

The column exists on \code{summary_metrics_v2} and the dropdown can
plot it.  The live v3 writer does not fill it.  On the current
database the field is null and the table prints \code{--}.
The working ATM rank on the page is the heatmap ATM column in
raw-IV mode, not this field.  Wing-spread leaves that column blank.


\section{The charts}


\subsection{Summary table}

One row per expiry \code{2nd}--\code{10th}.
\code{exp_date} is printed as \code{M/DD}.
Numeric cells are four decimals; \code{backbone_pct} is a percent or
\code{--}.

The table repaints every 2 seconds.  The fitted numbers are the
latest persisted snapshot, so they move on the writer clock
(30 seconds, and only after the live stream id changes), not on
every paint.  The 6-hour delta columns are the exception: they are
recomputed on each paint.

Read: this is the latest fitted snapshot, not a path.  Use it to see
the current ATM level, whether skew is call- or put-sided, and
whether the smile is curved up or down.

\subsection{Vol percentile heatmap}

Rows are expiries, nearest at the top (\code{autorange=reversed}).
Columns, left to right:

\begin{center}
\textsc{P100}, \textsc{P75}, \textsc{P50}, \textsc{P25},
\textsc{ATM}, \textsc{C25}, \textsc{C50}, \textsc{C75},
\textsc{C100}.
\end{center}

ATM uses moneyness bucket 0 (either type; first remaining row).
The other cells require matching type and bucket.
In wing-spread mode the ATM column is left blank (\code{n/a}):
that metric is \(\sigma-\sigma^{\mathrm{ATM}}\), so ranking the
ATM cell is either identically zero or a leftover of the
bucket-0 versus backbone-ATM mismatch.  Switch to raw IV to
see ATM level ranks.

Each cell is an empirical percentile of a scalar \(x\) against a
cached history list \(H\):

\begin{equation}
x
=
\begin{cases}
\sigma-\sigma^{\mathrm{ATM}}, & \text{wing-spread mode},\\
\sigma, & \text{raw-IV mode},
\end{cases}
\qquad
\mathrm{pctl}
=
100\cdot
\frac{\#\{h\in H:h<x\}}{|H|}.
\end{equation}

The implementation is \code{bisect_left}: values equal to \(x\) count
as not-less, so a print sitting on many identical history points
ranks at the bottom of that pile.

The color scale is blue at 0, white at 50, red at 100.
Blank cells are annotated \code{n/a}.

Read: red means rich versus the chosen history; blue means cheap.
Because the default metric is wing spread, a red C100 is ``that call
wing is high versus ATM relative to history,'' not ``absolute IV is
high.''  The middle ATM column is \code{n/a} in that mode.
Switch to raw IV when the question is the level, not the
smile shape.

\subsection{How the history list \(H\) is built}

The cache reads \code{heatmap_fact} and groups observations by

\begin{center}
(mode, expiry, moneyness, option type, 30-minute clock bucket,
event bucket, lookback window, regime filter, day-of-week bucket).
\end{center}

Rules:

\begin{itemize}
    \item 6-hour lookback uses no clock bucket and no weekday
    bucket.  1-week / 1-month / full use Eastern 30-minute clock
    buckets and a Mon--Wed versus Thu--Fri split.  Sunday evening
    Globex is counted with Monday.
    \item Regime thresholds are the p50 / p80 / p90 of that expiry's
    backbone inside the lookback.  Filters:
    all; exclude top 20\%; exclude top 10\%; keep bottom 50\%; keep
    top 20\%.
    \item A group is kept only if it has at least
    \(n_{\min}\) finite points (5 / 5 / 10 / 20 by lookback).
    If \(n\ge 4\), points outside \([Q_1-3\,\mathrm{IQR},
    Q_3+3\,\mathrm{IQR}]\) are dropped, then \(n_{\min}\) is
    required again.
    \item Event bucket \code{*} is the unconditional history.
    Time mode adds \code{t:...}; surprise mode adds \code{s:...}.
\end{itemize}

Lookup fallback, in order: exact key; then drop the event bucket to
\code{*}; then also drop the weekday bucket to \code{*}.
A cell that used a fallback is counted in the status line
(``\(n\) cells used fallback buckets'').

The cache refreshes every 900 seconds.  It stays empty until the
snapshot span is at least 3 days.

\subsection{Upcoming-events banner}

A yellow strip lists every remaining event day through a fixed
14-calendar-day window (the \code{2nd}--\code{10th} DTE span), not
whatever expiries the latest snapshot happens to contain, from
Yaman \code{events.sqlite} (US parents: coupon note/bond auctions,
scorer families, and High / Very High releases; bills, TIPS, FRN, and
Fed speeches are dropped).  Today's names include hours-ahead.
The line ends with ``percentiles may reflect event premium.''
The window is a fixed 14 calendar days.  It is not cut off at the
furthest expiry on the latest snapshot.

\subsection{Event premium / weight table}

One row per expiry.  Columns: Time Score, Surprise Score, the two
bucket labels, nearest parent, top weights, and the top three active
events.

The score is an additive kernel over events this expiry still owns.
An event counts only if it is still ahead and prints before this
expiry (12:00 CT on the expiry date).  After the print the
contribution is zero.  Same-calendar-day expiries use the
calibrated \code{1st} weight; later living expiries use that
ordinal's weight.  There is no 0.3 floor and no tier
\(\tau^{\mathrm{before}}\) floor.

\begin{equation}
s
=
\sum_{i:\,0<\Delta t_i,\; T_e\le T_{\mathrm{exp}}}
w_i\,e^{-\Delta t_i/\tau^{\mathrm{before}}_i}.
\end{equation}

\(w_i\) is the \code{1st} row when the expiry date equals the
event date, otherwise the row for this ordinal.

Buckets, identical for time and surprise
(\code{event_params.json}):

\begin{center}
\begin{tabular}{@{}lc@{}}
\toprule
Bucket & Score interval \\
\midrule
\code{no_event} & \([0,\,1.627)\) \\
\code{pre_event} & \([1.627,\,3.238)\) \\
\code{event_peak} & \([3.238,\,999)\) \\
\bottomrule
\end{tabular}
\end{center}

Displayed labels are \code{t:...} or \code{s:...}.
Parents with calibrated weights include ADP, AUCTION, CHALLENGER,
CLAIMS, CPI, DURABLES, FOMC, GDP, HOUSING, ISM, JOLTS, NFP, PCE,
PMI\_DETAIL, PPI, REGIONAL\_FED, RETAIL, SENTIMENT.

Read: a high time score means the calendar is close to a
heavily-weighted release for that expiry.  It does not mean the
surface has already moved.

\subsection{Event-distribution panel}

Opened by clicking an event-summary row.  Two charts:

\begin{itemize}
    \item ATM backbone histogram: blue unconditional sample versus
    orange event-conditioned sample, with a dashed red line at the
    current backbone.
    \item Surface boxplots of the same history used by the heatmap,
    plus a red diamond at the live value labeled with its
    percentile.  Purple boxes are fallback buckets.
    Wing-spread mode omits the ATM box for the same reason the
    heatmap ATM cell is blank.
\end{itemize}

The text block reports whether the current backbone would pass the
regime filter, the p50/p80/p90 thresholds, and the top event
contributions.

\subsection{Vol-decay panel}

Opened by clicking a heatmap cell.  For a chosen horizon
\(\{0.5,1,2,3,6\}\) hours it shows two historical probabilities that
ATM vol is lower after that horizon: event-conditioned (nearest
parent, 30-minute offset bin) and unconditional.

The footnote on the page is binding: those probabilities come from
ATM transitions (BentoBachelier four-year history) and are reused as
a proxy for every moneyness cell.  They are not a live model of that
wing.

\subsection{Metric timeseries}

One metric from
\code{backbone}, \code{skew_index}, \code{curve_index},
\code{rr25/50/75/100}, \code{fly25/50/75/100},
\code{backbone_roc_5m/15m/30m}, \code{backbone_pct},
\code{ts_slope}.
Default range is 1 hour.  Each selected expiry is its own line,
except \code{ts_slope}, which is a single 2nd-minus-5th series.

Read: this is the path.  The table is the latest point on it.

\subsection{Vol smile}

Latest snapshot only.  \(x\) is strike; \(y\) is
\code{implied_vol_daily_bp}.
If option type is ``Call + Put,'' the two IVs at the same strike are
averaged.  Default expiries are \code{2nd}--\code{5th}.

Read: this is the raw smile in price space.  It is not ranked.
A wing can look high on the smile and still be a mid-percentile on
the heatmap if that shape is typical for this clock and regime.


\section{Event conditioning, in words}

Unconditional ranks ask: ``is this print high versus recent prints
at this clock / weekday / vol regime?''

Time-only ranks ask the same question inside history that had a
similar calendar proximity score, without using the surprise.

Surprise-aware ranks use the same live kernel as time-only:
printed events are already zero, so there is no post-print
\(|z|\) tail on the score.  The label still exists on the page.

Conditioning never rewrites IV.  If a red wing turns white after
switching to event mode, the working idea is that the richness is
ordinary for this event, not that the quote changed.


\section{Parameter register}

\begin{longtable}{@{}p{0.34\linewidth}p{0.22\linewidth}p{0.36\linewidth}@{}}
\toprule
Name & Value & Role \\
\midrule
\endfirsthead
\toprule
Name & Value & Role \\
\midrule
\endhead
Dashboard port & 8055 & Local V6 page. \\
Tailscale & \code{:8443} \(\to\) 8055 & Remote URL. \\
Writer interval & 30 seconds & New table numbers.  Also
requires a new live stream id. \\
Adapter poll / stale & 1s / 10s & Live options feed. \\
Dashboard poll & 2 seconds & Repaint of the top table.
Values stay put until the next persist. \\
DB path & \code{C:/uikitxv2_cache/gamma_tracker/gamma_tracker_v3.db} &
Writer + read-only V6. \\
Live options host & \code{100.79.30.21:8001} & AccessAPI. \\
Price fallback & TY TBBO mid & If options \code{futures_mid} missing. \\
IV unit & daily BP, multiplier 1 & No annualization step. \\
Expiry labels & \code{2nd}--\code{10th} & Front remaining expiry is
\code{2nd}. \\
Delta targets & 0.50 / 0.375 / 0.25 / 0.125 / 0.05 & After binning. \\
RR / fly weights (raw) & 0.75 / 0.50 / 0.25 / 0.10 & Then /1.6. \\
ROC windows & 5 / 15 / 30 minutes & \code{2nd}--\code{5th} only. \\
\code{ts_slope} & \(B_{2\mathrm{nd}}-B_{5\mathrm{th}}\) & Status line
and chart. \\
Heatmap-fact sample & every 60 snapshots & Percentile history. \\
Raw retention & 30 days & Snapshots pruned; facts kept. \\
Tracker live / stale & 120s / 600s & Banner color. \\
Cache refresh & 900 seconds & Percentile rebuild. \\
Cache min history & 3 days & Else heatmap warming. \\
Default heatmap lookback & 1 month & TOD + DOW on. \\
Lookback \(n_{\min}\) & 5 / 5 / 10 / 20 & 6h / 1w / 1mo / all. \\
IQR fence & 3.0 & History outlier drop. \\
Regime thresholds & p50 / p80 / p90 of backbone & Per expiry and
lookback. \\
Event lookahead cap & 14 days & Banner and scorer. \\
Event calendar TTL & 1 hour & Scorer reload. \\
Event lookback & 12 hours & Calendar only; printed events score 0. \\
Structural future floor & 0.0 & Distant events decay to 0. \\
Tier \(\tau^{\mathrm{before}}\) floors & unused & Calibrated tau only. \\
Event-day weight & \code{1st} row & Same calendar date as the print. \\
Option expiry cutoff & 12:00 CT & Not 23:59 ET. \\
Score buckets & 1.627 / 3.238 & no\_event / pre\_event /
event\_peak. \\
Decay horizons & 0.5, 1, 2, 3, 6 hours & Click panel. \\
Event params date & 2026-03-25 & \code{event_params.json}. \\
\bottomrule
\end{longtable}


\section{What we are watching}

The page is for one question: is the live surface rich, cheap, or
ordinary versus its own history, and is that richness just the
calendar?

The types below are ways to read the same charts.  They are not
validated states and they are not long or short.

\subsection{Put-skew bid}

\(\mathrm{RR}_m<0\) on the table, and the left-hand put columns
(P75 / P100) are red on wing-spread while the matching calls are
not.  The working idea: downside vol is being bid versus upside.

\subsection{Call-skew bid}

The opposite: positive RR, red C75 / C100, quiet puts.
The working idea: upside vol is the expensive wing.

\subsection{Rich or cheap wings versus ATM}

Wing-spread mode, both wings red (or both blue).  The ATM column
is blank in this mode; confirm the ATM \emph{level} on raw IV.
The working idea: the smile is more (or less) curved than usual.
Confirm on \code{fly} / \code{curve_index} and on the smile chart.

\subsection{ATM vol shock}

\code{backbone} jumps on the timeseries and
\code{backbone_roc_5m} / \code{15m} is large, while the wing-spread
heatmap stays mostly white.  The working idea: the whole surface
repriced; the shape did not.  Raw-IV mode will go red together.

\subsection{Unexplained wing extension}

ATM is mid or quiet, but one far wing is stuck red and
\code{fly} is rising.  The working idea: that tail is being bid
without a matching ATM move.

\subsection{Event premium}

The yellow banner is populated, the \code{2nd} time-score is
\code{event_peak} or \code{pre_event}, and the unconditional heatmap
is red.  Switch to Event (time only).  If the reds fade, the working
idea: the richness is ordinary for this release.  If they stay red,
the surface is rich even among other event windows.

\subsection{Post-print surprise}

After the print the live score is 0, so time-only and surprise-aware
buckets both fall back to \code{no_event} unless another still-ahead
event is owned by that expiry.  There is no \(|z|\) afterglow.

\subsection{Term-structure disagreement}

\code{ts_slope} moves while a single expiry's backbone does not, or
the front heatmap row disagrees with \code{5th}--\code{7th}.
The working idea: the event or the shock is concentrated in one
tenor.

\subsection{Smile versus heatmap}

The smile looks steep in strike space, but the heatmap is mid.
The working idea: that steepness is typical for this clock / regime
/ event bucket.  Believe the heatmap for ``is this unusual?'';
believe the smile for ``what does the quote look like right now?''

\subsection{Decay panel versus a red cell}

A red ATM or wing with a high event-conditioned
\(P(\mathrm{vol\ lower})\) is a historical fade-vol pattern, not a
trade.  The panel is ATM history reused for every cell.  Treat it as
a reminder that rich event vol has often come in, not as an order.

\subsection{Quiet then ignition}

Banner \textsc{tracker: live}, scores \code{no_event}, heatmap mostly
white, then backbone ROC and one wing light up together.
The working idea: a new impulse in vol, not leftover event premium.
Early after a cache refresh this can also just be a thin bucket.

\subsection{Stale or warming}

\textsc{stale} / \textsc{offline}, or \textsc{cache: warming}.
Do not read ranks.  The quotes may still be on the table; the
history comparison is not ready.

These are reading habits.  The page never becomes a long or short.


\section{Ideas we are trying next}
\label{sec:ideas}

V6 tells us where a surface point sits versus history.  The working
question after that is more trader-shaped: for a given strike and
expiry, what are we paying for, what looks unusually priced, what is
likely to decay or reprice, and is the gamma we would buy plausibly
worth the theta and IV we pay?

That is a decision-support question, not a buy/sell painted onto the
heatmap.  The rest of this section writes the objects and the
experiments we want to try.  None of them are on the live page.
We will run them, look at the results, and rewrite this note.

The path we have in mind, if the ideas hold up, is
\[
\text{where is vol now?}
\;\to\;
\text{where should it be given the state?}
\;\to\;
\text{where is this node likely to go?}
\;\to\;
\text{is the gamma plus the mark change worth the price?}
\]
The first of those is V6.  The next three are what we are trying.


\section{Definitions}

These are working definitions for the experiments below.  If a
definition is awkward in the data, we change it.

\subsection{Canonical surface}

The live buckets do not match their targets.  Column ATM is
\(|\delta|\ge 0.60\), while backbone ATM is the print closest to
\(|\delta|=0.50\).  Bucket 25 is \(|\delta|\in[0.45,0.60)\) but the
kept strike aims at \(0.375\), and the same mismatch repeats down
the wing.  So a historical ``25-delta'' cell is not always the same
object.

The idea is to stop bucketing first.  Fit a smooth
\(\sigma_t(x,T)\) on each expiry, then evaluate it at fixed nodes
\[
|\Delta|\in\{0.50,\,0.375,\,0.25,\,0.125,\,0.05\}.
\]
Because the stored unit is already daily-bp normal vol, we will try a
Bachelier coordinate
\[
x=\frac{K-F_t}{\sigma_{\mathrm{ATM}}\sqrt{T}}
\]
rather than equity log-moneyness.  The screen can still talk in
delta.  The thing we store and compare is the interpolated node.

\subsection{Tenor}

``2nd'' at Monday 09:00 and ``2nd'' at Tuesday 15:45 are different
gamma/theta contracts.  The working tenor for any model is minutes
to expiry, or variance-time.  Ordinals stay as labels on the page.

\subsection{Absolute richness}

The V6 percentile is one absolute-richness screen:
\(\mathrm{pctl}(\sigma\mid\mathrm{hard\ bucket})\).
The idea we want to try is the residual against a conditional
expectation,
\[
R^{\mathrm{abs}}_{i,t}
=
\sigma_{i,t}
-
\widehat{\mathbb{E}}[\sigma_i\mid\mathrm{state}_t].
\]
State can include clock, DTE, event proximity, ATM regime, realized
vol, weekday, and the recent futures path.  The question is: is this
option's vol high or low given the environment?  We will try to
estimate that with partial pooling or a GAM, not by cutting the
current buckets finer.  Ten or twenty rows in a month are not twenty
independent samples.

\subsection{Local richness}

This is the strike-selection idea.  Fit today's smile while
excluding or downweighting strike \(i\), then
\[
R^{\mathrm{local}}_{i,t}
=
\sigma_{i,t}-\hat\sigma^{\mathrm{surface}}_{i,t},
\qquad
Z^{\mathrm{local}}_{i,t}
=
\frac{R^{\mathrm{local}}_{i,t}-\mathrm{median}(R^{\mathrm{local}}_i)}
{\mathrm{MAD}(R^{\mathrm{local}}_i)}.
\]
A wing can be cheap versus history and still locally expensive on
this smile.  Example: fitted 4.83\,bp, actual 5.08\,bp, local
residual \(+0.25\)\,bp.

\subsection{Variance time and forward variance}

\(\mathrm{ts\_slope}=B_{2\mathrm{nd}}-B_{5\mathrm{th}}\) is a level
gap.  For normal vol the more economic objects are
\[
V(T)=\sigma_N(T)^2 T,
\qquad
\mathrm{ForwardVar}_{T_1,T_2}
=
\frac{V(T_2)-V(T_1)}{T_2-T_1}.
\]
A high front IV is not automatically rich if that expiry owns an
event the next one does not.  What we want to know is how much
incremental variance the market is assigning to the gap.

\subsection{Event variance}

V6 already asks whether current IV is unusual for an event like
this.  The extra question is how many \(\mathrm{bp}^2\) the event is
contributing.  Write
\[
V(T)
=
v_{\mathrm{c}}T
+
\sum_{e:T_e<T}J_e^2.
\]
A neighbouring pair gives
\(J_e^2\approx V(T_+)-V(T_-)-v_{\mathrm{c}}\Delta T\).
With more expiries we can try
\[
V(T_j)
=
\alpha T_j
+
\sum_e \beta_e\,\mathbf{1}\{T_e<T_j\}
+\eps_j.
\]
Nine expiries and overlapping releases will alias.  We treat
\(\hat J_e^2\) as a readout to compare against its own history, not
as something tightly identified from one snapshot.

\subsection{Forward outcomes}

For each \((t,K,T)\) we want a later value of the same object.  Two
books, because they are not the same experiment:

\begin{itemize}
    \item same delta / same \(x\) at \(t+h\): smile research;
    \item the exact contract bought at \(t\): anything that looks
    like P\&L.
\end{itemize}

A rally turns today's 25-delta call into a later 50-delta call.
Working labels:
\begin{align}
Y^{\mathrm{IV}}_{i,t,h}
&=
\sigma_{i,t+h}-\sigma_{i,t},\\
Y^{\mathrm{spread}}_{i,t,h}
&=
(\sigma_i-\sigma_{\mathrm{ATM}})_{t+h}
-
(\sigma_i-\sigma_{\mathrm{ATM}})_t,\\
Y^{\mathrm{local}}_{i,t,h}
&=
R^{\mathrm{local}}_{i,t+h}-R^{\mathrm{local}}_{i,t}.
\end{align}
We will start with \(30\mathrm{m}\), \(1\mathrm{h}\), \(2\mathrm{h}\),
plus next-event / \(+5\mathrm{m}\) / \(+30\mathrm{m}\) after the
print and EOD, and add more horizons only if those are useful.

\subsection{Decay as a distribution}

The live decay panel is \(P(\mathrm{ATM\ vol\ lower})\) reused on
every wing.  The idea is a per-node forecast of \(\Delta\sigma_h\):
mean, median, \(P(\Delta\sigma<0)\), a left tail, and
\(Q_{0.1}/Q_{0.5}/Q_{0.9}\).

\subsection{Gamma carry and vol mark}

IV change alone is not the trade.  A 0.2\,bp drop can lose on vega
and still pay on gamma if TY realizes enough path.  Cheap IV can
still be a bad long-gamma if realized variance is cheaper.

For a locally delta-hedged option we will try the usual first
pieces,
\[
d\Pi
\approx
\Theta\,dt
+\tfrac12\Gamma(dF)^2
+\mathrm{Vega}\,d\sigma
+\cdots.
\]
Under the Bachelier local identity
\(\Theta\,dt\approx-\tfrac12\Gamma\sigma_N^2\,dt\), so
\[
\Pi^{\mathrm{gamma}}
\approx
\tfrac12\Gamma\bigl(QV_h-\sigma_N^2 h\bigr),
\qquad
QV_h=\sum_j(\Delta F_j)^2
\]
at the hedge clock we would actually use.  Endpoint
\(100\to 101\to 100\) is zero net move and still path.  Units of
\(QV_h\) and \(\sigma_N^2 h\) have to match; \(\Gamma\) is
normal-vol gamma.

Call
\[
G_i=\tfrac12\Gamma_i\bigl(\mathbb{E}[QV_h]-\sigma_i^2 h\bigr)
\]
the gamma-carry piece and
\[
M_i=\mathrm{Vega}_i\,\mathbb{E}[\Delta\sigma_i]
\]
the vol-mark piece.  Then, dropping vanna, volga, jumps, and
discrete-hedge error,
\[
\mathbb{E}[\Pi_i]\approx G_i+M_i-\mathrm{costs}.
\]
That is the composition we want to try later.  It is not a live V6
number.

\subsection{Four combined readings}

If both forecasts exist, the interesting cases are the usual four:
cheap with IV expected up and realized above implied; cheap with
both legs against (value trap); rich with both legs for a short;
rich while realized is still above implied (dangerous short).  Same
status as Section~8: ways to read the objects, not orders.


\section{Experiments}

Each item below is something we will try.  If it does not help, we
drop or rewrite it.

\subsection{E1. Interpolated nodes versus today's buckets}

Build \(\sigma_t(x,T)\) on each snapshot, evaluate at the fixed
delta nodes, and keep minutes-to-expiry on every row.  Also write
surface history denser than every 60th snapshot; the present
30-minute facts are thin for later labels.

What we look at: do the same nodes line up across days?  Does the
double-ATM problem go away?  How much does a node move when the
nearest listed strike jumps?

\subsection{E2. Conditional fair value}

Estimate \(\widehat{\mathbb{E}}[\sigma_i\mid\mathrm{state}]\) with a
GAM or another partial-pooling fit on clock, DTE, ATM, event score,
and weekday.  Compare \(R^{\mathrm{abs}}\) and its rank to the V6
hard-bucket percentile.

What we look at: does the residual still look like ``rich/cheap''
after the environment is taken out, and is the comparison set less
starved than a 10-row clock bucket?

\subsection{E3. Local residual as a strike screen}

Compute \(R^{\mathrm{local}}\) and \(Z^{\mathrm{local}}\) on each
snapshot.  See whether large positive local residuals later mean
revert on \(Y^{\mathrm{local}}\) more than a high V6 percentile
does.

What we look at: can two strikes in the same expiry have opposite
local residuals while sharing the same heatmap color?

\subsection{E4. Forward variance and event variance}

Replace the economic reading of \(\mathrm{ts\_slope}\) with \(V(T)\)
and \(\mathrm{ForwardVar}\).  Around CPI / NFP / FOMC, estimate
\(J_e^2\) from neighbouring expiries and from the short regression,
and compare today's \(J_e^2\) to its own history.

What we look at: does the front look ``rich in IV'' but ordinary in
event variance, or the reverse?  How unstable is \(\hat\beta_e\)
when two releases sit in the same gap?

\subsection{E5. Forward-label books}

Build the two books in the definitions: same-\(x\) and
same-contract.  Horizons as above.  This is the dataset everything
after this experiment reads.

What we look at: how often the two books disagree after a large TY
move, and whether \(Y^{\mathrm{local}}\) actually goes negative
after an extreme \(Z^{\mathrm{local}}\).

\subsection{E6. Decay distributions}

On that dataset, try in order: the current conditional-history
baseline; ridge; a GAM; quantile regression of
\(Q_{0.1},Q_{0.5},Q_{0.9}(\Delta\sigma\mid X)\).  Features we will
try first are surface state (level, spread, local residual, ATM,
variances, event variance), path (5/15/30 minute
\(\Delta\sigma\) and \(\Delta R^{\mathrm{local}}\), plus
\(\mathrm{ROC}_{5\mathrm{m}}-\mathrm{ROC}_{30\mathrm{m}}/6\)),
underlying (signed and absolute TY, realized variance, range, a
jump flag), event state, and contract state (DTE, greeks, bid/ask,
quote age).

Aggression / other order-flow is not in this first feature list.
Those studies already failed their own gates; joining them is a
later experiment if we still want it.

What we look at: does any simple model beat the existing ATM-proxy
panel on a walk-forward day or event split, in MAE / quantile
coverage / \(P(\Delta\sigma<0)\)?  Trees only if 0--3 lose.

\subsection{E7. Realized quadratic variation}

Forecast \(QV_h\) at the hedge clock we would actually use, not
end-to-end \(\Delta F\).  Same walk-forward rule as E6.

What we look at: how much of \(G_i\) is just a bad QV forecast
versus a real implied-versus-realized gap.

\subsection{E8. Combined expected hedged P\&L}

Add \(G_i+M_i-\mathrm{costs}\) on the same-contract book.  Show the
pieces, not one score.  Read them with the four combined cases
above.

What we look at: do the two pieces cancel in the cases we already
call value traps, and does \(\mathbb{E}[\Pi]\) after costs survive a
day-blocked walk-forward at all?

\subsection{E9. Event-day scenario reprice}

After E4--E8 have something to stand on, take historical
\((\Delta F,\,\Delta\mathrm{IV}_{\mathrm{ATM}},\,\Delta\mathrm{skew},\,\Delta\mathrm{curve})\)
around CPI / NFP / FOMC and reprice today's contract under those
paths.  Report \(P(\Pi>0)\), \(\mathbb{E}[\Pi]\), and a
10/90 band.  The event sample will be small.


\section{How we will run them}

Walk-forward by day or by event.  We will not randomly split
30-second or 30-minute rows.

In the research set we drop only bad data: crossed or stale marks,
failed IVs.  We keep economic tails.  The live \(3\times\mathrm{IQR}\)
fence can stay on V6 as a descriptive prototype; for these
experiments the extremes are part of the question.

Simple models have to beat the current conditional-history baseline
before we trust a more complicated one.
 
If an idea needs a new page, it is a later isolated copy, not an
edit of live V6.  Nothing here is a long or short.


\section{What we would want on a strike, if this works}

A click would show the pieces so we can see why a contract looks
interesting: IV; absolute and wing-spread ranks; local residual and
\(Z^{\mathrm{local}}\); conditional fair IV; the 30m and 2h
\(\Delta\mathrm{IV}\) distribution; expected QV versus the implied
variance budget; \(G\), \(M\), \(\mathbb{E}[\Pi]\) after costs;
\(P(\Pi>0)\); bid/ask and edge/cost; event-variance percentile;
confidence.  Ranking, if we rank at all, would use those pieces and
ratios such as \(\mathbb{E}[\Pi]/\mathrm{SD}(\Pi)\) and
\(\mathbb{E}[\Pi]/\mathrm{spread}\), not one collapsed score.

That card is a target if the experiments produce objects worth
showing.  It is not a spec for the current page.


\section{Order we will try them}

\begin{enumerate}
    \item Canonical interpolated surface; minutes-to-expiry; denser
    history (E1).
    \item Local residual and \(Z^{\mathrm{local}}\) (E3).
    \item Conditional fair value (E2).
    \item Total variance, forward variance, event variance (E4).
    \item Forward-label books (E5).
    \item Simple decay distributions (E6).
    \item QV forecast (E7).
    \item Combined \(G+M-\mathrm{costs}\) (E8).
    \item Event-scenario reprice (E9).
    \item Anything more complicated only if the simple versions
    help.
\end{enumerate}

The first cluster --- fair value, local residual, forward and event
variance --- is the one we care about before any decay model.  That
alone would already change how we pick a strike.  We will find out
whether it does.

\end{document}
